\documentclass[11pt,a4paper]{article}

% -----------------------------------------------------------------------------
% Required Packages
% -----------------------------------------------------------------------------
\usepackage[utf8]{inputenc}
\usepackage[margin=0.85in]{geometry}
\usepackage{graphicx}
\usepackage{xcolor}
\usepackage{listings}
\usepackage{booktabs}
\usepackage{tabularx}
\usepackage{fancyhdr}
\usepackage{hyperref}
\usepackage{tcolorbox}
\usepackage{caption}
\usepackage{titlesec}

% -----------------------------------------------------------------------------
% Color Palette & Custom Styling
% -----------------------------------------------------------------------------
\definecolor{collegeblue}{RGB}{30, 58, 138}
\definecolor{secondaryblue}{RGB}{30, 41, 59}
\definecolor{accentamber}{RGB}{245, 158, 11}
\definecolor{codebg}{RGB}{248, 250, 252}
\definecolor{codecomment}{RGB}{100, 116, 139}
\definecolor{codekeyword}{RGB}{219, 39, 119}
\definecolor{codestring}{RGB}{16, 185, 129}
\definecolor{lightborder}{RGB}{226, 232, 240}

\hypersetup{
    colorlinks=true,
    linkcolor=collegeblue,
    urlcolor=accentamber,
    citecolor=secondaryblue
}

% Section title styling
\titleformat{\section}
  {\normalfont\Large\bfseries\color{collegeblue}}{\thesection}{1em}{}[\titlerule]
\titleformat{\subsection}
  {\normalfont\large\bfseries\color{secondaryblue}}{\thesubsection}{1em}{}

% Header and Footer setup
\pagestyle{fancy}
\fancyhf{}
\fancyhead[L]{\small \textbf{Web Technology (2304362L)} --- Practical No. 4}
\fancyhead[R]{\small AY 2025--26}
\fancyfoot[L]{\small MIT Academy of Engineering, Pune}
\fancyfoot[R]{\small Page \thepage}
\renewcommand{\headrulewidth}{0.5pt}
\renewcommand{\footrulewidth}{0.5pt}

% -----------------------------------------------------------------------------
% Listings Setup (NodeJS / JavaScript)
% -----------------------------------------------------------------------------
\lstset{
    backgroundcolor=\color{codebg},
    basicstyle=\ttfamily\footnotesize,
    breaklines=true,
    captionpos=b,
    commentstyle=\color{codecomment}\itshape,
    frame=single,
    rulecolor=\color{lightborder},
    keywordstyle=\color{codekeyword}\bfseries,
    stringstyle=\color{codestring},
    numbers=left,
    numberstyle=\tiny\color{gray},
    numbersep=8pt,
    showstringspaces=false,
    tabsize=2,
    xleftmargin=12pt,
    framexleftmargin=8pt
}

% -----------------------------------------------------------------------------
% Document Body
% -----------------------------------------------------------------------------
\begin{document}

% -----------------------------------------------------------------------------
% TITLE / COVER PAGE
% -----------------------------------------------------------------------------
\begin{titlepage}
    \centering
    \vspace*{0.5cm}
    
    {\Large \textbf{MIT ACADEMY OF ENGINEERING, PUNE}}\\[0.2cm]
    {\large (An Autonomous Institute Affiliated to Savitribai Phule Pune University)}\\[0.2cm]
    {\textbf{School of Computer Engineering}}\\[1.2cm]
    
    \begin{tcolorbox}[colback=collegeblue,colframe=collegeblue,arc=3mm,width=\textwidth,halign=center]
        \vspace{0.3cm}
        {\LARGE \color{white} \textbf{LABORATORY REPORT}}\\[0.3cm]
        {\large \color{accentamber} \textbf{Course: Web Technology (2304362L) | T.Y. B.Tech | AY 2025--26}}
        \vspace{0.3cm}
    \end{tcolorbox}
    
    \vspace{1.2cm}
    
    {\LARGE \textbf{Practical No. 4}}\\[0.3cm]
    {\huge \bfseries \color{collegeblue} Server-Side Validation Application}\\[0.2cm]
    {\Large \textbf{using NodeJS Core Modules, Events, \& Streams}}\\[0.4cm]
    {\textit{Topic: Asynchronous Student Registration Server with Audit Stream Logging}}\\[1.5cm]

    \begin{minipage}{0.85\textwidth}
        \renewcommand{\arraystretch}{1.5}
        \begin{tabular}{|p{4.2cm}|p{8.5cm}|}
            \hline
            \textbf{Student Name} & Ayush \\ \hline
            \textbf{PRN / Roll No.} & 2304362 \\ \hline
            \textbf{Division / Batch} & T.Y. B.Tech Computer Engineering (Batch B) \\ \hline
            \textbf{Course Outcome} & CO3: Server-side Development using NodeJS \\ \hline
            \textbf{Lab Hours / Marks} & 8 Hours / 10 Marks \\ \hline
            \textbf{Date of Submission} & 27 / 10 / 2026 \\ \hline
            \textbf{Live Deployed URL} & \url{http://localhost:5000} \\ \hline
            \textbf{GitHub Repository} & \url{https://github.com/Ayush/wt-nodejs-server-validation} \\ \hline
        \end{tabular}
    \end{minipage}

    \vfill
    {\small \textbf{Academic Year 2025--26}}\\[0.2cm]
    {\footnotesize School of Computer Engineering, MIT Academy of Engineering, Alandi (D.), Pune --- 412105}
\end{titlepage}

\newpage

% -----------------------------------------------------------------------------
% SECTION 1: AIM & PROBLEM STATEMENT
% -----------------------------------------------------------------------------
\section{Aim and Problem Statement}

\subsection{Aim}
To design, develop, and deploy an asynchronous NodeJS HTTP server using native core modules (\texttt{http}, \texttt{fs}, \texttt{events}, \texttt{stream}, \texttt{path}, \texttt{os}) and modular custom modules that demonstrates strict server-side validation of incoming student registration payloads, proper HTTP status code handling (\texttt{200 OK}, \texttt{201 Created}, \texttt{400 Bad Request}, \texttt{409 Conflict}), event-driven audit logging, readable stream piping for data exports, and file-based JSON persistence.

\subsection{Problem Statement in Own Words}
Client-side form validation is inherently insecure as malicious agents or direct HTTP clients (such as Postman or cURL) can bypass frontend rules. Robust web systems require authoritative server-side validation. The core engineering goals accomplished are:
\begin{enumerate}
    \item \textbf{Native HTTP Server Without External Frameworks}: Building an HTTP web server using NodeJS core \texttt{http.createServer}, routing GET, POST, DELETE, and OPTIONS methods with CORS headers.
    \item \textbf{Authoritative Server-Side Validation Engine}: Enforcing 9 multi-tier validation checks (regex for names, 7--10 digit PRN, RFC email regex, 10-digit Indian mobile, enum checks for department and academic year, numeric range for CGPA, and non-empty skills arrays).
    \item \textbf{Proper HTTP Status Code Semantics}: Returning \texttt{201 Created} on success, \texttt{400 Bad Request} with a structured error array on validation failures, \texttt{409 Conflict} on duplicate PRN/Email, and \texttt{404 Not Found} on invalid endpoints.
    \item \textbf{NodeJS Event-Driven Architecture}: Extending \texttt{EventEmitter} to construct an asynchronous audit logger that emits domain events (\texttt{request}, \texttt{studentRegistered}, \texttt{validationFailed}, \texttt{duplicateAttempt}, \texttt{streamExport}).
    \item \textbf{NodeJS Streams Implementation}: Ingesting request payloads via readable streams, streaming CSV exports via \texttt{Readable.from().pipe(res)}, and streaming raw log files using \texttt{fs.createReadStream().pipe(res)}.
    \item \textbf{File System Persistence}: Asynchronously saving and retrieving records from \texttt{data/students.json} using \texttt{fs.promises}.
\end{enumerate}

% -----------------------------------------------------------------------------
% SECTION 2: SYSTEM ARCHITECTURE & DESIGN
% -----------------------------------------------------------------------------
\section{System Architecture and Design}

\subsection{Server Architecture and Request Lifecycle}
The architecture separates request handling, validation, storage, and event logging:

\begin{figure}[htbp]
\centering
\begin{tcolorbox}[colback=codebg,colframe=collegeblue,arc=2mm,title=\textbf{Figure 1: NodeJS Server Architecture and Stream Data Flow}]
\small
\begin{verbatim}
+-----------------------------------------------------------------------+
|                    Client HTTP Request (Browser / cURL)               |
|      POST /api/register  { fullName, prn, email, department, ... }    |
+-----------------------------------------------------------------------+
                                   |
                         (Readable Stream Chunks)
                                   v
+-----------------------------------------------------------------------+
|                      NodeJS HTTP Server (server.js)                   |
|  - Assembles chunks via req.on('data') and req.on('end')              |
|  - Routes to /api/register, /api/students, /api/stream/csv, etc.      |
+-----------------------------------------------------------------------+
                                   |
                                   v
+-----------------------------------------------------------------------+
|               Custom Validator Module (modules/validator.js)          |
|  - Regex pattern checks (Name, PRN, Email, Phone)                     |
|  - Range check (CGPA [0.00 - 10.00]) & Enum checks (Department, Year) |
|  - Input sanitization (XSS mitigation)                                |
+-----------------------------------------------------------------------+
                 | (Invalid)                          | (Valid)
                 v                                    v
+----------------------------------+   +--------------------------------+
|      HTTP 400 Bad Request        |   | Storage Manager (storage.js)   |
|  - Emits 'validationFailed' event|   | - Checks PRN/Email uniqueness  |
|  - Returns structured error JSON |   +--------------------------------+
+----------------------------------+          | (Duplicate)    | (Unique)
                                              v                v
                                      +---------------+  +--------------+
                                      | HTTP 409      |  | HTTP 201     |
                                      | Conflict      |  | Created      |
                                      +---------------+  +--------------+
                                                               |
                                                               v
+-----------------------------------------------------------------------+
|                 Event-Driven Audit Logger (modules/logger.js)         |
|  - Custom EventEmitter listening for 'studentRegistered' / 'stream'   |
|  - Appends timestamped audit log lines to data/server.log via fs      |
+-----------------------------------------------------------------------+
\end{verbatim}
\end{tcolorbox}
\end{figure}

\subsection{Technologies Used}
\begin{table}[htbp]
\centering
\small
\renewcommand{\arraystretch}{1.3}
\begin{tabularx}{\textwidth}{l p{3.2cm} X}
\toprule
\textbf{Technology} & \textbf{Type} & \textbf{Role in Practical No. 4} \\
\midrule
\textbf{NodeJS} & Runtime (v24.21.0) & Asynchronous non-blocking JavaScript execution engine. \\
\textbf{\texttt{http} Module} & Node Core & Creates native HTTP server, parses URLs, dispatches routes, sets status codes. \\
\textbf{\texttt{events} Module} & Node Core & Implements decoupled event-driven logging using \texttt{EventEmitter}. \\
\textbf{\texttt{stream} Module} & Node Core & Ingests chunked request bodies and streams CSV/log data via \texttt{pipe()}. \\
\textbf{\texttt{fs} / \texttt{fs.promises}} & Node Core & Asynchronous file I/O for JSON student storage and audit log persistence. \\
\textbf{\texttt{os} / \texttt{process}} & Node Core & Extracts CPU cores, memory RSS, and uptime metrics for \texttt{/api/stats}. \\
\bottomrule
\end{tabularx}
\caption{Technologies Used in Practical No. 4}
\end{table}

\newpage

% -----------------------------------------------------------------------------
% SECTION 3: KEY IMPLEMENTATION & CODE EXPLANATION
% -----------------------------------------------------------------------------
\section{Implementation and Key Code Explanation}

\subsection{Stream Request Body Ingestion (\texttt{server.js})}
Demonstrating NodeJS stream chunk handling and flood protection:

\begin{lstlisting}[language=JavaScript, caption={Stream Request Body Parser (server.js)}]
function parseRequestBody(req) {
  return new Promise((resolve, reject) => {
    let body = '';
    req.on('data', (chunk) => {
      body += chunk.toString();
      // 1MB flood protection limit
      if (body.length > 1e6) {
        req.connection.destroy();
        reject(new Error('Payload too large'));
      }
    });

    req.on('end', () => {
      if (!body.trim()) return resolve({});
      try {
        resolve(JSON.parse(body));
      } catch (err) {
        reject(new Error('Malformed JSON syntax in request body'));
      }
    });

    req.on('error', (err) => reject(err));
  });
}
\end{lstlisting}

\subsection{Custom Server-Side Validation Module (\texttt{modules/validator.js})}
Validating payload properties against strict academic criteria:

\begin{lstlisting}[language=JavaScript, caption={Validation Rules Engine (modules/validator.js)}]
function validateRegistration(data) {
  const errors = {};
  const sanitized = {};

  // Rule 1: Full Name (Alphabets and spaces only, 3-50 chars)
  const rawName = (data.fullName || '').trim();
  if (!rawName) errors.fullName = 'Full Name is required.';
  else if (!/^[A-Za-z\s]{3,50}$/.test(rawName))
    errors.fullName = 'Full Name must contain only alphabets and spaces (3-50 chars).';
  else sanitized.fullName = sanitize(rawName);

  // Rule 2: PRN (7-10 alphanumeric characters)
  const rawPRN = (data.prn || '').trim();
  if (!rawPRN) errors.prn = 'PRN / Roll Number is required.';
  else if (!/^[A-Za-z0-9]{7,10}$/.test(rawPRN))
    errors.prn = 'PRN must be 7 to 10 alphanumeric characters.';
  else sanitized.prn = sanitize(rawPRN).toUpperCase();

  // Rule 3: CGPA (Float range 0.00 to 10.00)
  const rawCGPA = parseFloat(data.cgpa);
  if (isNaN(rawCGPA) || rawCGPA < 0.0 || rawCGPA > 10.0)
    errors.cgpa = 'CGPA must be a valid numeric value between 0.00 and 10.00.';
  else sanitized.cgpa = parseFloat(rawCGPA.toFixed(2));

  return { isValid: Object.keys(errors).length === 0, errors, sanitizedData: sanitized };
}
\end{lstlisting}

\subsection{Event-Driven Audit Logger (\texttt{modules/logger.js})}
Extending \texttt{EventEmitter} to construct an asynchronous logging subsystem:

\begin{lstlisting}[language=JavaScript, caption={EventEmitter Audit Logger (modules/logger.js)}]
class ServerAuditEmitter extends EventEmitter {
  constructor() {
    super();
    this.setupListeners();
  }

  setupListeners() {
    this.on('studentRegistered', ({ studentId, prn, name, dept }) => {
      this.writeLog('SUCCESS', `Student Registered -> ID: ${studentId} | PRN: ${prn} | Name: "${name}"`);
    });

    this.on('validationFailed', ({ ip, errors, endpoint }) => {
      this.writeLog('WARN', `Validation Failed at ${endpoint} from IP ${ip} | Flagged: [${Object.keys(errors)}]`);
    });
  }

  writeLog(level, message) {
    const line = `[${new Date().toISOString()}] [${level}] ${message}\n`;
    fs.appendFile(LOG_FILE_PATH, line, (err) => {
      if (err) console.error('Failed writing to server.log:', err);
    });
  }
}
\end{lstlisting}

\newpage

% -----------------------------------------------------------------------------
% SECTION 4: TESTING & VERIFICATION MATRIX
% -----------------------------------------------------------------------------
\section{Testing and Test Cases}

\begin{table}[htbp]
\centering
\small
\renewcommand{\arraystretch}{1.3}
\begin{tabularx}{\textwidth}{l p{2.8cm} p{3.2cm} X c}
\toprule
\textbf{Test ID} & \textbf{Scenario} & \textbf{Request Payload} & \textbf{Expected HTTP Status \& Outcome} & \textbf{Status} \\
\midrule
\textbf{TC-01} & Valid Record Registration & Name: "Ayush", PRN: "2304362", CGPA: 9.42 & \texttt{201 Created}: Record persisted to \texttt{students.json}, audit event logged & \textbf{PASS} \\
\textbf{TC-02} & Multiple Validation Errors & Numbers in name, bad PRN, CGPA: 15.0 & \texttt{400 Bad Request}: Returns structured error dictionary for 9 invalid fields & \textbf{PASS} \\
\textbf{TC-03} & Duplicate PRN Conflict & Submitted existing PRN "2304362" & \texttt{409 Conflict}: "PRN '2304362' is already registered" & \textbf{PASS} \\
\textbf{TC-04} & Duplicate Email Conflict & Submitted existing email & \texttt{409 Conflict}: "Email already registered" & \textbf{PASS} \\
\textbf{TC-05} & CSV Stream Export & \texttt{GET /api/stream/csv} & \texttt{200 OK} with \texttt{Transfer-Encoding: chunked} streaming CSV file & \textbf{PASS} \\
\textbf{TC-06} & Server Log Stream & \texttt{GET /api/stream/logs} & \texttt{200 OK} streaming content from \texttt{data/server.log} via readable stream & \textbf{PASS} \\
\textbf{TC-07} & Runtime System Stats & \texttt{GET /api/stats} & \texttt{200 OK} returning Node version, uptime, memory RSS, and CPU count & \textbf{PASS} \\
\textbf{TC-08} & Record Deletion & \texttt{DELETE /api/students/:id} & \texttt{200 OK}: Record removed; or \texttt{404 Not Found} for unknown ID & \textbf{PASS} \\
\bottomrule
\end{tabularx}
\caption{System Test Execution Summary}
\end{table}

% -----------------------------------------------------------------------------
% SECTION 5: OUTPUT SCREENSHOTS
% -----------------------------------------------------------------------------
\section{Output Screenshots}

\begin{figure}[htbp]
\centering
\IfFileExists{screenshots/screenshot_valid_201.png}{%
    \includegraphics[width=0.75\textwidth]{screenshots/screenshot_valid_201.png}
}{%
    \begin{tcolorbox}[colback=codebg,colframe=lightborder,arc=2mm,height=5.2cm,valign=center,halign=center]
        \large \textbf{[ Screenshot 1: Valid Student Registration \& HTTP 201 Created Response ]}\\[0.3cm]
        \normalsize \textit{Attach screenshot showing successful POST request, green HTTP 201 Created badge, latency, formatted JSON payload, and populated records table}
    \end{tcolorbox}
}
\caption{Successful Student Registration Returning HTTP 201 Created and Persisted to Server}
\end{figure}

\begin{figure}[htbp]
\centering
\IfFileExists{screenshots/screenshot_invalid_400.png}{%
    \includegraphics[width=0.75\textwidth]{screenshots/screenshot_invalid_400.png}
}{%
    \begin{tcolorbox}[colback=codebg,colframe=lightborder,arc=2mm,height=5.2cm,valign=center,halign=center]
        \large \textbf{[ Screenshot 2: Server-Side Validation Failure \& HTTP 400 Bad Request Array ]}\\[0.3cm]
        \normalsize \textit{Attach screenshot showing Server Response Inspector with ruby-red HTTP 400 Bad Request, error count, and structured error breakdown}
    \end{tcolorbox}
}
\caption{Server Validation Engine Rejecting Invalid Payload with HTTP 400 Bad Request}
\end{figure}

\begin{figure}[htbp]
\centering
\IfFileExists{screenshots/screenshot_conflict_stream.png}{%
    \includegraphics[width=0.75\textwidth]{screenshots/screenshot_conflict_stream.png}
}{%
    \begin{tcolorbox}[colback=codebg,colframe=lightborder,arc=2mm,height=5cm,valign=center,halign=center]
        \large \textbf{[ Screenshot 3: Duplicate Record Conflict (HTTP 409) \& NodeJS Stream Export ]}\\[0.3cm]
        \normalsize \textit{Attach screenshot showing Server Response Inspector with HTTP 409 Conflict on duplicate PRN/Email, and active log stream viewer}
    \end{tcolorbox}
}
\caption{Duplicate Registration Rejected with HTTP 409 Conflict and Live Audit Log Stream}
\end{figure}

\newpage

% -----------------------------------------------------------------------------
% SECTION 6: CONCLUSION & LEARNINGS
% -----------------------------------------------------------------------------
\section{Testing and Conclusion}

\subsection{What Works}
\begin{itemize}
    \item \textbf{Native NodeJS HTTP Server}: Zero third-party runtime dependencies; built with core modules.
    \item \textbf{Authoritative Validation Engine}: Multi-rule checks protecting against bad, malicious, or malformed data.
    \item \textbf{REST Status Code Semantics}: Accurate use of \texttt{201 Created}, \texttt{400 Bad Request}, \texttt{409 Conflict}, \texttt{200 OK}, and \texttt{404 Not Found}.
    \item \textbf{Decoupled Event Architecture}: Subclassing \texttt{EventEmitter} for asynchronous non-blocking audit logging.
    \item \textbf{High-Throughput Streams}: Streaming CSV and log data directly via readable streams without memory overhead.
    \item \textbf{File System Persistence}: Reading and writing persistent JSON records with \texttt{fs.promises}.
\end{itemize}

\subsection{Limitations}
\begin{itemize}
    \item \textbf{File-Based Storage}: Flat JSON file storage is optimal for demonstration but requires relational databases (PostgreSQL/MySQL) or MongoDB for distributed enterprise workloads.
    \item \textbf{Payload Limit}: Body stream parser enforces a 1MB ceiling to guard against denial-of-service memory saturation.
\end{itemize}

\subsection{What I Learned}
\begin{itemize}
    \item Understood the NodeJS asynchronous event loop and stream-based HTTP request lifecycle.
    \item Mastered custom module design and the \texttt{EventEmitter} architectural pattern.
    \item Learned how to stream file data efficiently over HTTP with chunked transfer encoding.
    \item Gained firsthand appreciation of why server-side validation is indispensable for web security.
\end{itemize}

% -----------------------------------------------------------------------------
% SECTION 7: AI USE DECLARATION (Mandatory Table)
% -----------------------------------------------------------------------------
\section{AI Use Declaration}

\begin{table}[htbp]
\centering
\renewcommand{\arraystretch}{1.5}
\begin{tabularx}{\textwidth}{|p{4.5cm}|X|}
\hline
\textbf{AI tool and prompt used} & \textbf{Antigravity AI (Gemini 3.8)}\par
\textit{Prompt:} ``Design, develop and deploy an asynchronous NodeJS HTTP server demonstrating server-side validation, custom modules, file system persistence, EventEmitter audit logging, and stream piping.'' \\ \hline
\textbf{One output you checked and how} & Checked the chunked request body parser (\texttt{req.on('data')} and \texttt{req.on('end')}) by transmitting an incomplete JSON string via PowerShell to verify that the server properly catches syntax errors and returns HTTP 400 Bad Request instead of crashing the process. \\ \hline
\textbf{One thing you corrected or rejected, and why} & Corrected the duplicate check mechanism so that PRN uniqueness comparisons are case-insensitive (\texttt{toUpperCase()}) and Email comparisons are lowercase-normalized (\texttt{toLowerCase()}), preventing duplicate registrations with differing letter cases. \\ \hline
\end{tabularx}
\caption{Mandatory AI Use Declaration Table}
\end{table}

\end{document}
